\section{Frobenius comparison and smooth nonvanishing}
\label{fr:section}

The scalar and two-slot estimates lead to a contradiction through a
comparison in positive characteristic. We first prove that certain
fixed divisors, jets, and a movable-curve witness cannot coexist.
This theorem and its proof are independent of nonvanishing, abundance,
minimal models, and logarithmic subadditivity. We then construct its
inputs from the counterexample geometry and obtain smooth canonical
nonvanishing.

\subsection{The finite-data Frobenius theorem}
\label{common:frobenius-tools}

We now work with arbitrary fixed data, independent of the scaling
models and divisors of Section~\ref{sec:jets}. The comparison concerns
two orders of vanishing of an actual determinant. Fixed ordinary jets
will give a nonzero determinant after reduction. A small polarization bounds the
rank on the diagonal and therefore forces a large order there.
A fixed movable curve on a blowup bounds every section in each
determinant factor, producing the opposite estimate.

For a vector bundle on a smooth projective curve, its slope is
degree divided by rank. The least and greatest Harder--Narasimhan
slopes are denoted by $\mu_{\min}$ and $\mu_{\max}$.
For a line bundle on a smooth variety, to generate jets through
order $a$ at a point means surjectivity to its stalk modulo
the $(a+1)$-st power of the maximal ideal.

\begin{theorem}[Finite-data Frobenius incompatibility]
\label{common:finite-data-frobenius}
Let $W$ be a smooth connected projective complex variety of
dimension $n\ge2$. Let $L,H$ be rational Cartier divisors, with
$H$ ample, let $S$ be an integral Cartier divisor, and fix an
integer $q>0$ and real numbers $r>1$, $\epsilon>0$. Suppose
\begin{equation}\label{common:intersection-bounds}
 \begin{gathered}
 H^n\le(2\epsilon)^n,\qquad
 K_WH^{n-1}\le2H^n/r,\qquad
 (2L+qS)H^{n-1}\le4H^n,\\
 (14\epsilon)^n<\tfrac14,\qquad 80\epsilon<\tfrac34.
 \end{gathered}
\end{equation}
Assume that the following fixed data exist.
\begin{enumerate}[label=\textup{(\roman*)}]
\item A smooth integral complete-intersection flag from $W$ to
a curve $C$, whose successive divisor classes are $h_i=l_iH$
for fixed positive integers $l_i$, with
$\mu_{\min}(\Omega_W^1|_C)\ge0$.
\item On the projective bundle
\[
 Z=\mathbf P\bigl(\mathcal O_W(S)_1\oplus
                         \mathcal O_W(S)_2\bigr)\longrightarrow W\times W,
 \qquad \xi=c_1(\mathcal O_Z(1)),\qquad M=L_1+L_2+q\xi,
\]
use the convention that the direct image of $\mathcal O_Z(a\xi)$
is the $a$-th symmetric power of the displayed sum. A smooth
point $z_*$ in the complement of the two axes and finitely many
sections of $k_0M$, for some integer $k_0>0$ with $k_0L$ Cartier,
generate jets through order $(2n+2)k_0$ at $z_*$.
\item A point $x\in W$, its blowup
$\pi_x:\widehat W\to W$ with exceptional divisor $J$, and a
curve class
\[
 \gamma=f_*(A_1\cdots A_{n-1}),
\]
where $f:V\to\widehat W$ is one fixed birational morphism,
$V$ is smooth integral projective, and the $A_i$ are ample
integral Cartier divisors, such that $J\cdot\gamma>0$ and
\begin{equation}\label{common:movable-input}
 \pi_x^*(L+K_W/2+\lambda S)\cdot\gamma
 \le40\epsilon J\cdot\gamma\qquad(0\le\lambda\le q).
\end{equation}
It suffices to check the two endpoints in this affine inequality.
\end{enumerate}
These data cannot coexist.
\end{theorem}

There is no relation required between $x$ and the two base
coordinates of $z_*$. The divisors $L,S,K_W$ are not assumed
nef, and $K_W$ is not assumed pseudo-effective. The theorem
starts with the actual flag and movable witness, rather than
with geometric conditions from which one might construct them.
All data in its statement are fixed before the residual
characteristic tends to infinity.

\subsubsection{Spreading the fixed witnesses}

We prove the theorem in the next five steps. Put
$\Lambda=\prod_{i=1}^{n-1}l_i$. The flag satisfies
\begin{equation}\label{common:flag-degrees}
 [C]=\Lambda H^{n-1},\qquad h_i\cdot C=l_i\Lambda H^n.
\end{equation}
Choose one sufficiently ample integral Cartier divisor $T_0$ so
that every $T_0+aS$, $0\le a\le(k_0-1)q$, has a section
nonvanishing at each of the two base coordinates of $z_*$.
Fix these finitely many sections. All models, morphisms, divisors,
points, the flag, the complete-intersection witness for $\gamma$,
and the jet and filler sections spread over an integral finitely
generated $\mathbf Z$-algebra of characteristic zero. Shrink its
spectrum so that the fibers and flag are smooth projective and
geometrically integral, the birational morphisms remain birational,
the fixed ample bundles remain ample, and the finite jet
surjection and nonvanishings persist. Birationality is preserved
by spreading an isomorphism between dense opens. All displayed
intersection numbers are constant.

We also preserve $\mu_{\min}(\Omega_W^1|_C)\ge0$.
One way is to spread its characteristic-zero Harder--Narasimhan
filtration, make all its factors locally free, and use openness
of their semistability on curves; their degrees and ranks are
fixed. Equivalently, a single relative ample twist globally
generates the curve bundle. Every quotient of rank $a$ then has
degree bounded below by $-ad$ for a fixed integer $d$.
A negative-degree quotient has one of finitely many ranks and
degrees. The proper relative Quot schemes for those Hilbert
polynomials have images missing the generic point: after removal
of torsion, a negative-degree coherent quotient there would
give a negative-degree vector-bundle quotient. Removing those
finitely many images proves the same openness assertion directly.

This base has closed points in arbitrarily large characteristics.
Take algebraic closures of their residue fields and exclude
denominator primes. Retain the notation for the reductions.
The curve $\gamma$ continues to pair nonnegatively with every
effective divisor: pull that divisor back under the fixed
birational morphism and intersect with its fixed ample divisors.
This tests effective divisors on the reduction itself. No
specialization of a characteristic-zero effective cone is used.

\begin{lemma}\label{common:positive-characteristic-jets}
For a sufficiently large residual characteristic $p$, set
\[
 N_p=\lfloor p/k_0\rfloor,\qquad B_p=k_0N_pL+T_0.
\]
The sections of $pq\xi+B_{p,1}+B_{p,2}$ generate jets through order
$(2n+2)k_0N_p$ at $z_*$.  In particular, they surject onto the quotient
of its local ring by the $p$-th powers of all regular parameters.
\end{lemma}

\begin{proof}
Multiply $N_p$ copies of the fixed jet system.  Every monomial of
total degree at most $(2n+2)k_0N_p$ is a product of $N_p$ monomials
of degree at most $(2n+2)k_0$.  Taking sections with those leading
monomials gives a triangular system, ordered by total degree.
Starting at degree zero and removing the error in each successive
degree proves the required jet surjection. This argument takes
place in the local ring and never divides by factorials.

Put $d_p=(p-k_0N_p)q$, so $0\le d_p\le(k_0-1)q$.
Among the sections fixed before reduction, choose
\[
 u_0\in H^0(W,\cO_W(T_0)),\qquad
 v_{d_p}\in H^0(W,\cO_W(T_0+d_pS))
\]
nonzero at the first and second base coordinates of $z_*$,
respectively. Their tensor product belongs to the summand of
first-slot weight zero in the projective-bundle formula; its
fiber monomial is the $d_p$-th power of the second coordinate.
It defines a section of
$d_p\xi+T_{0,1}+T_{0,2}$ nonzero at $z_*$, since $z_*$ lies
off both axes. Multiplying the product jet system by this section
changes its divisor from $N_pk_0M$ to
$pq\xi+B_{p,1}+B_{p,2}$. Multiplication is invertible on the
local jet algebra. Only the finitely many previously fixed
values of $d_p$ occur.

Since $\dim Z=2n+1$, the
quotient by the $p$-th powers of regular parameters has top total
degree $(2n+1)(p-1)$.  For large $p$ this is at most
$(2n+2)k_0N_p$. The underlying ordinary/Frobenius ideal comparison is
also recorded in \cite[proof of Proposition~2.12, Equation~(2.8)]{MustataSchwede2014}.
\end{proof}

\subsubsection{Full rank off the diagonal}

Let $F:W\to W'$ be relative Frobenius to the base-field twist.
It is finite flat because $W$ is smooth.  Write $S'$ for the twist
of $S$, so $F^*S'=pS$, and define
\[
 \mathcal E=F_*\cO_W(B_p),\quad s=\rk\mathcal E=p^n,\quad R=s^2,
 \qquad U_a=H^0(W,\cO_W(B_p+paS)).
\]
Projection formula gives an evaluation map on $W'\times W'$:
\begin{equation}\label{common:two-slot-evaluation}
 \bigoplus_{\substack{a+b=q\\a,b\ge0}}
 U_a\otimes U_b\otimes
 \cO(-aS'_1-bS'_2)
 \longrightarrow \mathcal E\boxtimes\mathcal E.
\end{equation}

\begin{lemma}\label{common:generic-rank}
The map in Equation~\eqref{common:two-slot-evaluation} has generic rank $R$.
\end{lemma}

\begin{proof}
Trivialize the two summands defining $Z$ near the selected base
points.  Let $z$ be the torus coordinate, whose value at $z_*$ is
$z_0\ne0$.  By the projective-bundle formula, the sections in
Lemma~\ref{common:positive-characteristic-jets} are weight polynomials in $z$
with weight-$j$ coefficients in
\[
 H^0(W,B_p+jS)\otimes
 H^0(W,B_p+(pq-j)S),\qquad 0\le j\le pq.
\]
Let $A_*$ be the tensor product of the local algebras of the two
base Frobenius fibers.  The local quotient in
Lemma~\ref{common:positive-characteristic-jets} is
\[
 A_*[z]/(z^p-z_0^p),
\]
free over $A_*$ with basis $1,z,\ldots,z^{p-1}$.  Project to the
coefficient of $1$ in this basis.  Exactly the weights $j=pa$
survive, with multipliers $z_0^{pa}$.  Their coefficients are the
values of the corresponding summands of
Equation~\eqref{common:two-slot-evaluation}, in the chosen frames and the induced
twisted frames.  As the full jet map is surjective, these coefficients
span $A_*$, of dimension $p^{2n}=R$.  Thus the evaluation has full
rank at this pair of base points, and hence generically.  The
calculation uses a basis over the possibly nonreduced algebra $A_*$,
so reducedness of the Frobenius fibers is unnecessary.
\end{proof}

The jet calculation at $z_*$ proves generic full rank; the diagonal
estimate below will be tested at $(x',x')$, obtained from the
separately fixed point $x$.

\subsubsection{The diagonal filtration}

The Frobenius fiber product has the cartesian square
\begin{equation}\label{common:frobenius-square}
\begin{tikzcd}
 \Delta_F=W\times_{W'}W \arrow[r,"\mathrm{pr}_2"] \arrow[d,"\mathrm{pr}_1"'] & W \arrow[d,"F"]\\
 W \arrow[r,"F"'] & W'.
\end{tikzcd}
\end{equation}
Its reduced subscheme is the diagonal $W$.  Put
\[
 \mathcal L_p=
 \cO_{\Delta_F}(B_{p,1}+B_{p,2}+pqS_1).
\]
Write $h=F\circ\mathrm{pr}_1=F\circ\mathrm{pr}_2:
\Delta_F\to W'$. On this fiber product,
\[
 \cO(pS_1)\simeq h^*\cO_{W'}(S')\simeq\cO(pS_2).
\]
Consequently every pair of weights $a+b=q$ gives the same line bundle:
\[
 \cO_{\Delta_F}(B_{p,1}+B_{p,2}+paS_1+pbS_2)
 \simeq\mathcal L_p.
\]
On the diagonal of $W'\times W'$, every source twist of
Equation~\eqref{common:two-slot-evaluation} becomes $-qS'$.
Finite base change and projection formula identify the twisted target as
\[
 (\mathcal E\otimes\mathcal E)(qS')\simeq h_*\mathcal L_p.
\]
The product sections defining the columns restrict to global sections
of this one bundle $\mathcal L_p$. Their values at a diagonal point
therefore lie in the image of $H^0(\Delta_F,\mathcal L_p)$ in the
corresponding fiber of $h_*\mathcal L_p$. This proves that the rank
at every diagonal point is at most
\begin{equation}\label{common:diagonal-rank-upper}
 h^0(\Delta_F,\mathcal L_p).
\end{equation}

For a rank-$n$ bundle $V$ in characteristic $p$, denote by $A_j(V)$
the degree-$j$ part of its symmetric algebra modulo the $p$-th powers
of local generators.  This construction is independent of frame:
the $p$-th power of a linear combination is the sum of the $p$-th
powers in characteristic $p$.  Each $A_j(V)$ is locally free and a
quotient of $V^{\otimes j}$.  Set
\[
 u=n(p-1),\qquad e_j=\rk A_j(V),\qquad
 \sum_{j=0}^u e_j=p^n.
\]
Filtering $\mathcal L_p$ by powers of the ideal of the reduced
diagonal in $\Delta_F$ gives the graded bundles
\begin{equation}\label{common:diagonal-graded}
 \mathcal G_j=\cO_W(2B_p+pqS)\otimes A_j(\Omega_W^1),
 \qquad 0\le j\le u.
\end{equation}
Indeed, in smooth local coordinates the algebra is generated by
parameter differences with their $p$-th powers zero, and its
degree-one conormal is $\Omega_W^1$.  This coordinate calculation
can be made étale locally or in completions and identifies the
global associated graded.  This is the diagonal-ideal form of the canonical Frobenius filtration
\cite[Corollary~3.5]{KitadaiSumihiro2008}. The same calculation, using
only the bundle from the second factor, filters $F^*\mathcal E$ with pieces
\cite[Theorem~3.7]{Sun2008}
\begin{equation}\label{common:Frob-pullback-graded}
 \cO_W(B_p)\otimes A_j(\Omega_W^1).
\end{equation}

Complementary multiplication is a perfect pairing, giving
\begin{equation}\label{common:socle-pairing}
 A_{u-j}(V)\simeq A_j(V)^\vee\otimes(\det V)^{p-1}.
\end{equation}
On monomials, each exponent vector pairs with its complement to
$(p-1,\ldots,p-1)$.  The top line transforms by the character
$(\det V)^{p-1}$: this is immediate on diagonal matrices and hence
 identifies the character of $\mathrm{GL}_n$ on that line.

\subsubsection{One slope bound and one flag polynomial}

We bound the sections of every graded bundle in
Equation~\eqref{common:diagonal-graded} by its rank times the same scalar
polynomial in $p$. This uniformity lets us sum the ranks of the
graded pieces without an additional factor for their number.
We begin with a slope bound on the fixed curve; here slope means
ordinary degree divided by rank.

\begin{lemma}\label{common:truncated-slope}
There is a constant $c\ge0$, independent of the sufficiently large
residual characteristic $p$ and of $j$, such that for
$V=\Omega_W^1|_C$ one has
\[
 \mu_{\max}(A_j(V))\le(p-1)K_W\cdot C+nc.
\]
\end{lemma}

\begin{proof}
Write $F_C$ for absolute Frobenius of $C$.  Langer's instability
estimate \cite[Corollary~2.5]{Langer2004} gives
\begin{equation}\label{common:Langer-min}
 L_{\min}(V):=
 \lim_{e\to\infty}p^{-e}\mu_{\min}(F_C^{e*}V)
 \ge-\frac{c}{p-1}.
\end{equation}
To check the uniform constant, the discrepancy from
$\mu_{\min}(V)\ge0$ is at most
$(\rk V-1)\deg A_C/(p-1)$, where $A_C$ is a nef line bundle such
that $T_C\otimes A_C$ is globally generated.  Its degree can be
bounded in terms of the fixed genus, using a sufficiently positive
line bundle.  Here the rank is $n$, irrespective of the fact that the base
is a curve.

We also need, for every integer $i\ge0$,
\begin{equation}\label{common:tensor-min}
 \mu_{\min}(V^{\otimes i})\ge iL_{\min}(V).
\end{equation}
Fix $p$ and $i$.  Choose a line bundle $A_e$ of degree
$-\mu_{\min}(F_C^{e*}V)+O(1)$, rounded up with a fixed
genus-dependent additive constant, so that
$F_C^{e*}V\otimes A_e$ is globally generated.  This follows from
Serre duality: after subtracting any point, its minimum slope can
be made greater than $2g(C)-2$, which annihilates $H^1$ and makes
evaluation at that point surjective.  For every vector-bundle
quotient $Q$ of $V^{\otimes i}$, the bundle
$F_C^{e*}Q\otimes A_e^{\otimes i}$ is then globally generated.
Consequently
\[
 p^e\mu(Q)+i\deg A_e\ge0.
\]
Divide by $p^e$ and let $e\to\infty$ to prove
Equation~\eqref{common:tensor-min}.  This limit is taken for each fixed $p,i$;
there is no interchange of Frobenius-iteration and characteristic
limits.

Since $A_{u-j}(V)$ is a quotient of $V^{\otimes(u-j)}$, the perfect
pairing Equation~\eqref{common:socle-pairing} yields
\begin{align*}
 \mu_{\max}(A_j(V))
 &= (p-1)\deg\det V-\mu_{\min}(A_{u-j}(V))\\
 &\le(p-1)K_W\cdot C+\frac{(u-j)c}{p-1}\\
 &\le(p-1)K_W\cdot C+nc.
\end{align*}
\end{proof}

The divisors $B_p-pL=T_0-(p-k_0N_p)L$ range over a fixed finite
list.  Combining Equation~\eqref{common:intersection-bounds}, Equation~\eqref{common:flag-degrees}, Lemma~\ref{common:truncated-slope}
therefore gives, uniformly for $0\le j\le u$,
\begin{align}
 \mu_{\max}(\mathcal G_j|_C)
 &\le p(4+2/r)\Lambda H^n+O(1)\notag\\
 &\le M_p:=7p\Lambda H^n+c_0,\label{common:common-slope-bound}
\end{align}
where $c_0\ge0$ is fixed.  The coefficient $7$ provides harmless
room in this common bound.

\begin{lemma}\label{common:diagonal-sections}
For all sufficiently large $p$,
\[
 h^0(\Delta_F,\mathcal L_p)
 \le R\bigl(7^nH^n+O(1/p)\bigr)<R/4.
\]
The error constant is independent of the filtration index $j$.
\end{lemma}

\begin{proof}
Write the fixed flag as
\[
 W=W_n\supset W_{n-1}\supset\cdots\supset W_1=C,
 \qquad W_{n-k}\in|h_k|_{W_{n-k+1}}
 \quad(1\le k\le n-1).
\]
At step $k$, fix the preceding nonnegative integers
$b_1,\ldots,b_{k-1}$ and put
\[
 \mathcal F_k=
 \mathcal G_j|_{W_{n-k+1}}
 \Bigl(-\sum_{i<k}b_i h_i\Bigr).
\]
For each $b_k\ge0$, the divisor restriction sequence is
\[
 0\longrightarrow\mathcal F_k(-(b_k+1)h_k)
 \longrightarrow\mathcal F_k(-b_kh_k)
 \longrightarrow
 \mathcal F_k(-b_kh_k)|_{W_{n-k}}
 \longrightarrow0.
\]
Iterating it for $b_k=0,1,\ldots$ and then proceeding to the next
member of the flag bounds the section space by
\begin{equation}\label{common:flag-lattice-sum}
 h^0(W,\mathcal G_j)\le
 \sum_{b_1,\ldots,b_{n-1}\ge0}
 h^0\left(C,\mathcal G_j|_C\Bigl(-\sum_i b_i h_i|_C\Bigr)\right).
\end{equation}
At each stage the remainder is zero after sufficiently negative
ample twisting.  Its stopping index need not be uniform in $j,p$
or in the preceding twists: enlarging the nonnegative finite sums
to the displayed lattice sum preserves the inequality.

A summand in Equation~\eqref{common:flag-lattice-sum} vanishes when
$\sum_i b_i(h_i\cdot C)>M_p$, because then its maximum slope is
negative.  Each nonzero summand has at most $e_j(M_p+1)$ sections.
Indeed, evaluation at $\lfloor M_p\rfloor+1$ distinct points is injective;
the kernel has negative maximum slope after subtracting those
points.  By Equation~\eqref{common:flag-degrees} it follows that
\begin{align}
 h^0(W,\mathcal G_j)
 &\le e_j(M_p+1)
       \prod_{i=1}^{n-1}
       \left(1+\frac{M_p}{l_i\Lambda H^n}\right)\notag\\
 &=e_jp^n\bigl(7^nH^n+O(1/p)\bigr).
 \label{common:common-flag-polynomial}
\end{align}
The leading coefficient follows from
$\Lambda=\prod_i l_i$.  More importantly, the whole scalar
polynomial on the first line is the same for every $j$.

Sum over the filtration in Equation~\eqref{common:diagonal-graded}.  The ranks
satisfy $\sum_j e_j=p^n$, so Equation~\eqref{common:common-flag-polynomial} gives
the asserted bound with $R=p^{2n}$.  There is no extra factor for
the number of graded pieces.  Finally
\[
 7^nH^n\le(14\epsilon)^n<1/4
\]
by Equation~\eqref{common:intersection-bounds}, proving the strict inequality
for all large $p$.
\end{proof}

\subsubsection{The determinant has incompatible orders}

Choose $R$ generically independent columns of
Equation~\eqref{common:two-slot-evaluation} from bases of its source vector spaces.
Their determinant is a nonzero section
\[
 0\ne\sigma\in H^0(W'\times W',\mathcal Q_1\boxtimes\mathcal Q_2)
\]
of an actual external-product line bundle. To specify its factors,
let $m_i$ be the sum of the weights of the $R$ chosen columns in
slot $i$. The target determinant is
$(\det\mathcal E)^s\boxtimes(\det\mathcal E)^s$, and the
chosen source determinant is
$\cO(-m_1S')\boxtimes\cO(-m_2S')$. Thus
\[
 \mathcal Q_i=(\det\mathcal E)^{\otimes s}
                 \otimes\cO_{W'}(m_iS'),\qquad
 \lambda_i=m_i/R.
\]
Every column weight lies between $0$ and $q$, so
$0\le\lambda_i\le q$ independently of the columns chosen.

Under the numerical identification with the base-field twist,
\begin{equation}\label{common:slot-determinant-class}
 \frac{c_1(\mathcal Q_i)}R
 =\frac{c_1(\mathcal E)}s+\lambda_iS
 =\frac{B_p}p+\frac{p-1}{2p}K_W+\lambda_iS,
 \qquad 0\le\lambda_i\le q.
\end{equation}
The Frobenius first-Chern-class identity is also
\cite[Lemma~4.2]{Sun2008}; it is an identity of rational divisor classes.  For the second equality, use
Equation~\eqref{common:Frob-pullback-graded}.  Complementary degrees in
Equation~\eqref{common:socle-pairing} have total first Chern class
$e_j(p-1)K_W$.  Summing gives
\[
 \sum_jc_1(A_j(\Omega_W^1))=\frac{s(p-1)}2K_W.
\]
Thus $c_1(F^*\mathcal E)=sB_p+s(p-1)K_W/2$, and Frobenius pullback of
twisted divisor classes multiplies by $p$, as required.

For $0\le\lambda\le q$, the fixed movable inequality gives
\[
 (L+K_W/2+\lambda S)\cdot\gamma\le40\epsilon J\cdot\gamma,
\]
where pullback by $\pi_x$ is understood. This concerns only
intersections of fixed divisors and the fixed complete-intersection
witness, so it remains true on the reduction. The class in
Equation~\eqref{common:slot-determinant-class} differs from its
left-hand divisor at $\lambda=\lambda_i$ by
\[
 \frac1p\bigl(B_p-pL-K_W/2\bigr).
\]
Its numerator belongs to the fixed finite list
$T_0-aL-K_W/2$, $0\le a<k_0$. The resulting error after
intersection with $\gamma$ is therefore $O(p^{-1})$, uniformly
over the selected determinant columns.

Let $x'$ be the twisted point and let $0\ne\tau_i\in
H^0(W',\mathcal Q_i)$.  The strict transform of its effective divisor on
the blowup at $x'$ pairs nonnegatively with the twist of $\gamma$.
Since $J\cdot\gamma>0$, it follows that
\begin{equation}\label{common:slot-order}
 \frac{\ord_{x'}(\tau_i)}R\le40\epsilon+O(1/p).
\end{equation}
The constant is uniform in the columns and in $\tau_i$.  This
argument tests sections on the reduction itself; it does not
require them to lift to characteristic zero.

K\"unneth identifies the space containing $\sigma$ with
$H^0(W',\mathcal Q_1)\otimes H^0(W',\mathcal Q_2)$.  Choose bases adapted to the
orders of vanishing at $x'$ in each factor.  In each degree their
leading terms are linearly independent.  Tensor products of these
leading terms remain independent in each bidegree of the two
disjoint sets of parameters; terms of different bidegrees cannot
cancel.  Every nonzero tensor consequently has order at most the
sum of the two maximum basis orders.  Applying
Equation~\eqref{common:slot-order}, we obtain
\begin{equation}\label{common:determinant-order-upper}
 \ord_{(x',x')}\sigma\le R\bigl(80\epsilon+O(1/p)\bigr).
\end{equation}

On the other hand, Equation~\eqref{common:diagonal-rank-upper}, Lemma~\ref{common:diagonal-sections}
bound the matrix rank at $(x',x')$ by a number strictly smaller
than $R/4$.  Trivialize its line bundles near that point.  Invertible
constant row operations make more than $3R/4$ rows vanish at the
point.  Each entry of those rows belongs to the maximal ideal, so
every term of the determinant has order at least $3R/4$.  Therefore
\begin{equation}\label{common:determinant-order-lower}
 \ord_{(x',x')}\sigma\ge3R/4.
\end{equation}
Equations \eqref{common:determinant-order-upper} and
\eqref{common:determinant-order-lower} contradict
$80\epsilon<3/4$ for sufficiently large $p$.


This contradiction proves Theorem~\ref{common:finite-data-frobenius}.
The full-rank calculation used $z_*$, while the diagonal estimate
holds at every point and was tested at the separately chosen $x$.
The errors were controlled by one fixed finite list of divisor
classes and one fixed flag polynomial. Thus their constants do
not depend on $p$, the filtration index, or the chosen columns.

\subsection{The geometric witnesses and smooth nonvanishing}
\label{fr:geometric-witnesses}

We return to the notation and the hypothetical counterexample of
Section~\ref{sec:jets}. Theorem~\ref{common:finite-data-frobenius}
asks for fixed divisors, a small polarization and flag, finite two-slot
jets, and one actual strongly movable curve. We construct all these
inputs in characteristic zero. The positive parts on the scaling
models, rather than the original non-nef classes, provide the small
polarization. No minimal-model statement or pseudo-effective cone is
specialized to positive characteristic.

\begin{theorem}[The smooth nonvanishing step]
\label{thm:smooth-nonvanishing}
Assume the lower-dimensional real-boundary good-model hypothesis of
Assumption~\ref{mmp:lower-models}.
If \(X\) is a smooth projective complex variety of dimension \(n\) and
\(K_X\) is pseudo-effective, then \(\kappa(X,K_X)\geq0\).
\end{theorem}

In dimension zero the assertion is immediate.  On a smooth curve,
pseudo-effectivity of \(K_X\) gives \(g(X)\geq1\), hence
\(h^0(X,K_X)=g(X)>0\). We therefore suppose \(n\geq2\) and argue by
contradiction. We use the late terminal model \(Y\), the divisors
\(K=K_Y\) and \(A=A_Y\), and the scaling models of the preceding
section. Fix \(m>0\) with \(mK\) Cartier, and write
\[
 \mu_t=\vol(K_X+tA_X)^{1/n},\qquad
 L=r(K+tA).
\]
As before, \(r\) is a positive integer chosen so that
\(r\mu_t\longrightarrow\epsilon\) as \(t\downarrow0\).
Choose the rational number \(\epsilon>0\) sufficiently small for
Propositions~\ref{prop:scalar-jets} and~\ref{prop:large-seshadri}, and also
so that
\begin{equation}
\label{fr:epsilon}
 (14\epsilon)^n<\frac14,\qquad
 80\epsilon<\frac34.
\end{equation}
The strict inequalities leave room for the errors that tend to zero
with the characteristic.  We next fix a sufficiently large integer
\(q\) as in Proposition~\ref{prop:large-seshadri}, and then choose
a positive rational \(t\) small enough for the estimates below and for
\(r>qm+1\). Once those estimates hold, \(t\) and \(r\) remain fixed
through all the subsequent characteristic-zero constructions and
reductions.

\subsubsection{A small polarization and a cotangent flag}

A small ample polarization and a complete-intersection flag will
control the number of sections on the Frobenius thickening of the
diagonal. We construct them from the positive part on a scaling model.

Choose a smooth common resolution \(W\) of \(Y\) and the scaling model
for \(K+tA\). In this application \(K,A,L\) also denote their pullbacks to
\(W\), and
\[
 K_W=K+E,\qquad E\geq0.
\]
Let \(H_0\) be the pullback of the nef semiample positive part \(N_0\) of
\(L\) on its scaling model. The comparison of canonical pullbacks and
moving parts gives
\begin{equation}
\label{fr:limit-intersections}
 H_0^n\longrightarrow\epsilon^n,\qquad
 LH_0^{n-1}=H_0^n,\qquad
 K_WH_0^{n-1}=KH_0^{n-1}\leq\frac{H_0^n}{r}.
\end{equation}
To justify these intersections, recall that the maps
\(Y\dashrightarrow Y_t\) are small by the choice of the late model.
Every divisor exceptional over \(Y\) is therefore also exceptional
over \(Y_t\).  The differences discarded when passing to the moving
parts are exceptional over the scaling model, and the transform of
\(A\) is effective up to rational linear equivalence.  Intersecting
with \(H_0^{n-1}\) gives the displayed equalities and inequality.
In particular, since \(qm/r<1\),
\[
 (2L+qmK)H_0^{n-1}
 \le(2+qm/r)H_0^n<3H_0^n.
\]
Choose \(t\) small enough that also \(H_0^n<(2\epsilon)^n\), and
now fix \(t,r\) and the resolution \(W\). For a fixed ample divisor
\(H_{\mathrm{amp}}\), put \(H=H_0+\eta H_{\mathrm{amp}}\), where
\(\eta>0\) is rational. The strict margins just obtained, together
with \(K_WH_0^{n-1}\le H_0^n/r\), allow us to choose \(\eta\)
small enough that continuity gives
\begin{equation}
\label{fr:intersection-bounds}
 H^n\leq(2\epsilon)^n,\qquad
 K_WH^{n-1}\leq\frac{2H^n}{r},\qquad
 (2L+qmK)H^{n-1}\leq4H^n.
\end{equation}

We next choose the flag for the section estimates.  The canonical
divisor of \(W\) is pseudo-effective.  Generic
semipositivity, applied with empty boundary
\cite[Theorem~2.1]{CampanaPaun2015}, and restriction to a sufficiently
general complete-intersection curve give the following fixed data.
Apply restriction to the Harder--Narasimhan filtration and its factors
\cite[Theorem~6.1 and Remark~6.2]{MehtaRamanathan1982}. Choose successively
large divisible integers $l_i$ such that $h_i=l_iH$ are integral very ample
classes, and choose a smooth integral complete-intersection flag
\[
 W=W_n\supset W_{n-1}\supset\cdots\supset W_1=C,
 \qquad W_{i-1}\in |h_{n-i+1}|_{W_i},
\]
for which
\begin{equation}
\label{fr:generic-nef}
 \mu_{\min}(\Omega_W^1|_C)\geq0.
\end{equation}
Here and below slopes on \(C\) mean degree divided by rank. One may
obtain the flag by applying restriction to the finitely many
Harder--Narasimhan quotients of \(\Omega_W^1\).

\subsubsection{A movable witness for all determinant weights}

For the vanishing-order estimate, choose a very general point
\(x\in W\), away from the exceptional loci, and let
\(\pi_x:\widehat W\to W\) be its blowup, with exceptional
divisor \(J\). Set
\[
 D^+=2r(K+2tA)+2E.
\]
This big divisor has the scalar sections controlled by
Proposition~\ref{prop:scalar-jets}: pushing to \(Y\) identifies the
section spaces, since the added exceptional part does not change them.
If \(\rho=\vol(2r(K+2tA))^{1/n}\), then
\[
 \rho\leq4r\mu_t\leq8\epsilon
\]
for our sufficiently small \(t\). Thus every section of a sufficiently
divisible multiple of \(D^+\) has normalized order at \(x\) at most
\(4\rho\leq32\epsilon\).

We convert this order bound into a curve inequality.
It follows that \(\pi_x^*D^+-40\epsilon J\) is not pseudo-effective.
Otherwise its convex combination with the big class \(\pi_x^*D^+\)
would make \(\pi_x^*D^+-cJ\) big for some rational
\(32\epsilon<c<40\epsilon\), giving a section of excessive order.
Movable-curve duality \cite[Theorem~2.2]{BDPP2013} therefore supplies an
actual covering curve class \(\gamma\) on \(\widehat W\) such that
\begin{equation}
\label{fr:curve-test}
 (\pi_x^*D^+)\cdot\gamma
       <40\epsilon\,J\cdot\gamma,\qquad J\cdot\gamma>0.
\end{equation}
The strictness lets us fix an algebraic family that will survive
reduction.  More precisely, the negative pairing can first be detected
by a strongly movable curve: take the pushforward of a general complete
intersection of very ample divisors on one fixed smooth birational
model of \(\widehat W\). Fix this model, its divisors, and the resulting
covering family.  Thus the choice consists of finite algebraic data,
rather than a limiting real movable class.  Positivity of
\(J\cdot\gamma\) follows from the strict inequality and
pseudo-effectivity of \(\pi_x^*D^+\).

To turn this curve into the witness required by
Theorem~\ref{common:finite-data-frobenius}, set $S=mK$ on $W$;
it is an integral Cartier divisor. We verify all its weight inequalities
before proceeding to the jet construction. Put
$Q_\lambda=L+K_W/2+\lambda S$ for $0\le\lambda\le q$.
The equality $K_W=K+E$ gives
\begin{equation}\label{fr:slot-domination}
 D^+-Q_\lambda
 =(r-\tfrac12-\lambda m)K+3rtA+\tfrac32E.
\end{equation}
This class is pseudo-effective in characteristic zero: $r>qm+1$,
$K$ is pseudo-effective, and $A,E$ are effective up to the stated
equivalences. Pairing with the fixed strongly movable class therefore gives
\[
 \pi_x^*Q_\lambda\cdot\gamma
 \le \pi_x^*D^+\cdot\gamma
 <40\epsilon J\cdot\gamma,
 \qquad J\cdot\gamma>0.
\]
These are the required movable inequalities. Only the two endpoint
inequalities need to be retained as finite data, since the pairing is
affine in $\lambda$. The common theorem spreads these fixed numerical
inequalities and the actual birational complete-intersection witness;
it does not spread the pseudo-effective assertion used to obtain them.

\subsubsection{A finite two-slot jet system}

We must similarly realize the two-slot jet estimate by finitely many
sections before reducing.  Their products will then supply the jet
orders needed for arbitrarily large characteristics.

On \(W\times W\), let
\[
 Z=\PP(mK_1\oplus mK_2),\qquad
 \xi=\OO_Z(1),\qquad
 M=L_1+L_2+q\xi,
\]
with the symmetric-power convention for the projective bundle.
Proposition~\ref{prop:large-seshadri} gives a semiample big model whose
Seshadri constant exceeds \(2n+2\) at a very general torus point.
On the open set where its birational contraction is an isomorphism,
the jet interpretation of the Seshadri constant consequently gives
an integer \(k_0>0\) such that \(|k_0M|\) generates jets of order
at least \((2n+2)k_0\) at a fixed smooth torus point \(z_*\in Z\).
Indeed, on the ample model choose a rational number \(c\) strictly
between \(2n+2\) and its Seshadri constant. On the blowup of the
selected smooth point, the pullback of the ample class minus \(c\)
times the exceptional divisor is ample. Serre vanishing and the
exceptional-divisor sequence then give jets of order \(kc-1\) for
all sufficiently divisible large \(k\). Pulling sections back on the
isomorphic open gives the asserted bound.
Enlarge \(k_0\) to clear every denominator, in particular that of \(L\),
and fix finitely many sections realizing this jet surjection.


\subsubsection{Applying the comparison}

All the preceding choices are now fixed. With $S=mK$, the dimension,
intersection and flag hypotheses of
Theorem~\ref{common:finite-data-frobenius} are
Equations~\eqref{fr:epsilon}, \eqref{fr:intersection-bounds},
and~\eqref{fr:generic-nef}. The curve construction gives its movable
witness and both endpoint inequalities, and the finite jet system
gives the required surjection at $z_*$. This point need not lie over
the separately chosen testing point $x$. The proof of the common
theorem also fixes its auxiliary divisor $T_0$ and finitely many
filler sections before choosing any residual characteristic.

The common theorem now rules out these fixed data. Its contradiction
compares two orders of the same nonzero determinant. After reduction,
the Frobenius direct image \(\mathcal E\) used there has rank \(p^n\),
so \(\mathcal E\boxtimes\mathcal E\) has rank \(R=p^{2n}\). The fixed
jets give full evaluation rank \(R\), while the small polarization
bounds the rank on the diagonal by \(R/4\). Thus a nonzero maximal
minor vanishes to order at least \(3R/4\) at \((x',x')\), where \(x'\)
is the Frobenius image of the testing point \(x\). Its line bundle
is an actual external product. Applying the fixed curve inequality
to each factor bounds the same order above by
\(R(80\epsilon+O(p^{-1}))\), a contradiction for large \(p\).
The theorem proves these determinant estimates for the fixed witnesses
we have now supplied; it requires no additional geometric input here.

All constants are fixed before the residual prime varies.
The contradiction excludes the assumed counterexample and proves
Theorem~\ref{thm:smooth-nonvanishing}.
